\section{The universal minimum-degree theorem}
\label{sec:minimal-degree}

This section proves the part of \cref{thm:main} that quantifies over every
finite extension \(L/\Q\).  The exact \(U_1\) calculation enters only after
the complete source theorem has reduced the possibilities to two branches.

\begin{lemma}[Intersection fixed field]
\label{lem:intersection-fixed-field}
Let \(L/\Q\) be finite, put \(G_L=\Gal(\Qbar/L)\), and let
\begin{equation}
 H_L=\im\!\left(G_L\longrightarrow\Gal(K/\Q)\right).
\label{eq:def-HL}
\end{equation}
Then
\[
 K^{H_L}=K\cap L,\qquad
 [K\cap L:\Q]=[\WE:H_L].
\]
\end{lemma}

\begin{proof}
Because \(K/\Q\) is finite Galois, restriction identifies the image of
\(G_L\) with \(\Gal(K/K\cap L)\).  Its fixed field in \(K\) is therefore
\(K\cap L\).  The degree formula follows from the Galois correspondence and
\cref{eq:frozen-line-field}.
\end{proof}

\begin{theorem}[Sharp degree threshold]
\label{thm:degree}
If \(L/\Q\) is finite and \(\bigl(\Brq{L}\bigr)[2]\ne0\), then
\begin{equation}
 36\mid[K\cap L:\Q]\mid[L:\Q].
\label{eq:degree-chain}
\end{equation}
If \([L:\Q]=36\) and \(\bigl(\Brq{L}\bigr)[2]\ne0\), then
\begin{equation}
 L=K\cap L=K^{wU_1w^{-1}}
\label{eq:equality-field}
\end{equation}
for some \(w\in\WE\), with a unique double-six embedded in the fixed
algebraic closure.  Conversely, each field \(K^{wU_1w^{-1}}\) has degree
\(36\), and its Brauer quotient is nonzero.
\end{theorem}

\begin{proof}
Since \(Y_{\Qbar}\) is rational,
\(\Br(Y_{\Qbar})=0\).  The inflation and number-field Hochschild--Serre
identifications in \cref{lem:inflation,eq:brauer-picard-bridge}, applied with
\(k=L\), turn the hypothesis into
\[
 H^1(H_L,\Lambda)[2]\ne0.
\]
The complete 2-primary classification
\cite{SwinnertonDyer1993,ElsenhansJahnel2010Brauer} gives exactly the two
containments in \cref{tab:two-branches}, for a suitable \(w\in\WE\).

\begin{table}[htbp]
\centering
\begin{tabular}{@{}ccc@{}}
\toprule
2-primary quotient & forced containment & ambient index\\
\midrule
\(\Z/2\) & \(H_L\subseteq wU_1w^{-1}\) & \([\WE:U_1]=36\)\\
\((\Z/2)^2\) & \(H_L\subseteq wU_3w^{-1}\) & \([\WE:U_3]=720=36\cdot20\)\\
\bottomrule
\end{tabular}
\caption{Both complete classification branches force divisibility by \(36\),
but only the \(U_1\) branch can attain index \(36\).}
\label{tab:two-branches}
\end{table}

In the first branch, the subgroup index factors as
\begin{equation}
 [\WE:H_L]
 =36[U_1:w^{-1}H_Lw].
\label{eq:u1-index-factor}
\end{equation}
In the second,
\begin{equation}
 [\WE:H_L]
 =720[U_3:w^{-1}H_Lw].
\label{eq:u3-index-factor}
\end{equation}
Thus \(36\mid[\WE:H_L]\) in either case.  By
\cref{lem:intersection-fixed-field}, this is the first divisibility in
\cref{eq:degree-chain}.  The inclusion \(K\cap L\subseteq L\) and the tower
law give the second.

Now suppose, in addition, that \([L:\Q]=36\).  Every factor in
\cref{eq:degree-chain} is one,
so \(L=K\cap L\) and \([\WE:H_L]=36\).  The index-\(720\) branch is
impossible.  \Cref{eq:u1-index-factor} forces
\(H_L=wU_1w^{-1}\), and the fixed-field statement in
\cref{eq:equality-field} follows.

The exact configuration action verifies that \(U_1\) is self-normalizing and
is the stabilizer of exactly one double-six.  Therefore its conjugates are
indexed by
\[
 [\WE:N_{\WE}(U_1)]=[\WE:U_1]=36,
\]
and they correspond to the thirty-six double-sixes.  Distinct conjugate
subgroups have distinct fixed fields by Galois correspondence.  This proves
uniqueness of the embedded field attached to \(D\).  The fields are permuted
transitively by \(\WE=\Gal(K/\Q)\), hence form a single
\(\Q\)-isomorphism type.

Conversely, \(K^{wU_1w^{-1}}\) has degree \(36\).  The exact integral
calculation \(H^1(U_1,\Lambda)\cong\Z/2\), together with the arithmetic bridge
proved in \cref{sec:brauer-jump}, shows that its Brauer quotient is nonzero.
\end{proof}

\begin{corollary}[Sharp threshold]
\label{cor:sharp-threshold}
The smallest degree of a finite extension \(L/\Q\) for which
\(\bigl(\Brq{L}\bigr)[2]\ne0\) is \(36\).  Every such degree is divisible by \(36\).
\end{corollary}

\begin{proof}
The divisibility is \cref{eq:degree-chain}; existence in degree \(36\) is the
converse part of \cref{thm:degree}.
\end{proof}

\begin{remark}[Why one double-six is insufficient]
The existence of a selected \(U_1\)-fixed configuration proves attainment,
not universal minimality.  Without the \((\Z/2)^2\)--\(U_3\) row in
\cref{tab:two-branches}, the proof would leave an uncontrolled family of
Galois images.  The fact that \(720\) is itself divisible by \(36\) is what
allows the two branches to yield one sharp divisibility theorem.
\end{remark}
